% H9-T20: moving tame Kummer covers of positive genus.
% Run JAC1964-20261121; /root/h9_height; GPT-6 Astra Pro.

\subsection*{H9-T20. The same exact height norm for generic positive-genus Kummer covers}

\begin{theorem}[H9-T20]
Keep the abelian variety and Lefschetz pencil of H9-T19, over
an algebraically closed field $k$, with $\infty\notin\Sigma$.
Let $N\ge2$ be prime to the characteristic. Choose a positive
multiple $B=Nq_0$ of $N$ and distinct points
$b_1,\ldots,b_B\in\mathbf A^1(k)-\Sigma$, and set
\[
                     a(t)=\prod_{i=1}^{B}(t-b_i).
\]
Let $s$ be transcendental over $k$, put $F=\overline{k(s)}$,
and let $C/F$ be the smooth projective curve whose function
field is
\[
 L=F(t,u),\qquad u^N=(t-s)a(t),
                  \qquad K=F(t).
                                                               \tag{H9.101}
\]
Then $[L:K]=N$, the genus of $C$ is $B(N-1)/2$, and the cover
is unramified above $\Sigma$. On the full homologically
trivial rational codimension-$m$ group of $Y_L$, all cycles
are admissible and
\[
 h_L(z,w)=\frac1N h_K(\operatorname{cor}z,
                            \operatorname{cor}w),\qquad
 \operatorname{rad}h_L
       =\operatorname{cor}^{-1}(\operatorname{rad}h_K).
                                                               \tag{H9.102}
\]
In particular, for an abelian fourfold, accepted032 and
H9-T8 give full-group semipositivity over this function field
of positive genus. The same conclusions hold for the
geometric generic cyclic cover with the same branch
exponents and freely varying branch points. In
characteristic different from two, taking $N=2$ and $B=2g$
gives the geometric generic hyperelliptic cover of every
genus $g\ge1$.
\end{theorem}

\begin{proof}
We use the cohomological and Chow arguments of H9-T19, but
must now treat the constant ambient odd cohomology as well as
the variable pencil system. This extra term is not discarded.

\paragraph{The curve family is smooth after explicit tame normalization.}
Let
\[
 S=\mathbf A^1_s-(\Sigma\cup\{b_1,\ldots,b_B\}).
\]
On $\mathbf P^1_{[T:Z]}\times S$ consider the cyclic root
cover for the section
\[
 (T-sZ)\prod_{i=1}^{B}(T-b_iZ)\,Z^{N-1}
            \in H^0\bigl(\mathcal O(B+N)\bigr).
                                                               \tag{H9.103}
\]
Here $B+N=N(q_0+1)$, so the root line bundle is
$\mathcal O(q_0+1)$. Normalize this finite cover and denote
the result by $\mathscr C\to\mathbf P^1_t\times S$.
We verify smoothness over $S$ rather than assuming it for
the unnormalized equation.

At each finite branch section the equation, after an \'etale
extension adjoining an $N$th root of a unit, is
$v^N=t-b_i$ or $v^N=t-s$. Its normalization is smooth over
$S$. At infinity, write $z=1/t$ and rescale $u$ by
$z^{q_0+1}$. The local equation becomes
$w^N=z^{N-1}c(z,s)$ with $c$ a unit. An \'etale extension
removes the unit because $N$ is invertible. The normalization
of $w^N=z^{N-1}$ is explicitly
\[
                        z=v^N,\qquad w=v^{N-1},
\]
which is again smooth over $S$. Away from the branch sections
the cover is \'etale. These local descriptions prove that
$\mathscr C\to S$ is a smooth proper curve family and that
its finite map to $\mathbf P^1_t\times S$ is flat of degree
$N$. For flatness one may also use the same regular-target,
Cohen--Macaulay-source finite-morphism criterion as in
H9-T19a. The normalization is finite here, and the displayed
local normalizations verify its formation in this family.

The zero at the moving section has valuation one. Thus the
Kummer extension has degree $N$, for example by Eisenstein
at $t-s$, and all fibres are geometrically connected.
Every one of the $B+2$ branch points, including infinity,
is totally and tamely ramified: the exponent at infinity
is $N-1$, prime to $N$. Riemann--Hurwitz gives
\[
 2g(C)-2=-2N+(B+2)(N-1),
                 \qquad g(C)=\frac{B(N-1)}2.
\]
The branch sections avoid $\Sigma$ by construction.

Set
$\mathscr X=\widetilde A\times_{\mathbf P^1_t}\mathscr C$,
and write $p:\mathscr X\to S$ for its structure map.
As in H9-T19, it is smooth and proper over $S$: over the old
critical values the curve cover is \'etale, and elsewhere
the old pencil morphism is smooth. The finite map
$Q:\mathscr X\to\widetilde A\times S$ is flat of degree
$N$ and has deck group $\mu_N$.

\paragraph{The entire nontrivial deck part is computed.}
Fix $\ell$ invertible in $k$ and use
$\overline{\mathbf Q}_\ell$ coefficients for characters.
Put
\[
 H=H^{n-1}(A,\overline{\mathbf Q}_\ell),\qquad
 V=(R^{n-1}f_*\overline{\mathbf Q}_\ell)_{\mathrm{var}},
 \qquad r=\operatorname{rank}V>0.
\]
The condition-A and Leray arguments of H9-T19 apply verbatim
to $f_s:\mathcal X_s\to C_s$. The two outer Leray terms
are still deck invariant: the relevant sheaves are constant,
$H^0(C_s)$ is trivial, and every deck automorphism acts by
degree one on $H^2(C_s)$. The middle term is still permanent
from $E_2$. Unlike the rational-cover case,
$H\otimes H^1(C_s)$ may be nonzero and is retained below.

For every nontrivial deck character $\chi$, let
\[
                       M_\chi=\mathcal L_\chi(a(t)).
\]
Up to the harmless choice of inverse character, its piece
of the covering direct image is
$M_\chi\otimes\mathcal L_\chi(s-t)$.
At every old critical point the cover is unramified; at its
branch points the old pencil sheaves are lisse. Thus the
same stalk calculation as in H9-T19 gives the precise middle
extensions. The relative Leray identification and Katz's
parabolic recipe \S5.1.7(3)--(4) yield the lisse-sheaf
decomposition
\[
 (R^n p_*\overline{\mathbf Q}_\ell)_\chi
       =\bigl(H\otimes MC_\chi(M_\chi)\bigr)
                 \oplus MC_\chi(V\otimes M_\chi)
                         \quad\text{on }S.                 \tag{H9.104}
\]
The constant vector space $H$ carries trivial parameter
monodromy. Equation (H9.104) includes both the fixed and the
variable contributions and excludes extra punctual terms
by condition~A, not by a claim of generic cohomological
splitting for singular fibres.

\paragraph{Both terms have no finite-orbit vector.}
The rank-one middle extension $M_\chi$ is nonconstant and
irreducible. It has exactly $B\ge2$ nontrivial finite
singularities, each with nontrivial character $\chi$, and
is unramified at infinity because $N\mid B$. It is therefore
an object of Katz's category $\mathcal T_\ell$. By
Theorem~2.9.7 and \S5.1.4, $MC_\chi(M_\chi)$ is irreducible.
Applying the exact infinity rule (H9.97) with original rank
one and sum of finite drops $B$ gives
\[
 \operatorname{rank}MC_\chi(M_\chi)=B,\qquad
 MC_\chi(M_\chi)_\infty
           =J(\chi,2)\oplus J(\chi,1)^{\oplus(B-2)}.
                                                               \tag{H9.105}
\]
A finite power of inertia is nontrivial unipotent. Hence this
irreducible system has infinite image and no finite-orbit
vector, by the elementary argument in H9-T19. Tensoring it
with the constant space $H$ is a direct sum of copies and
still has no finite-orbit vector.

The sheaf $V\otimes M_\chi$ is also irreducible: tensoring
with a rank-one local system preserves irreducibility.
It is tame and unramified at infinity. At each old critical
point its inertia is the original rank-one transvection,
with drop one; at each $b_i$ its inertia is the nontrivial
scalar $\chi$, with drop $r$. With $q=|\Sigma|$, the sum
of finite drops is $q+Br$. Katz's same rule therefore gives
\[
 \begin{split}
 \operatorname{rank}MC_\chi(V\otimes M_\chi)&=q+Br,\\
 MC_\chi(V\otimes M_\chi)_\infty
   &=J(\chi,2)^{\oplus r}
        \oplus J(\chi,1)^{\oplus(q+(B-2)r)}.
 \end{split}                                                  \tag{H9.106}
\]
It too is irreducible with infinite image, and has no
finite-orbit vector. The conventions at infinity are the
same as in (H9.98); replacing $\chi$ by $\chi^{-1}$ does
not change any conclusion. No all-prime sign comparison or
orthogonal-monodromy classification is needed.

Equation (H9.104) now proves that the entire nontrivial deck
part has no finite-orbit vector. Every algebraic class on
$\mathcal X_F$ is fixed after a finite extension of $k(s)$.
The same field-of-definition argument as in H9-T19 gives
\[
 \operatorname{cl}\operatorname{CH}^m(\mathcal X_F)_{\mathbf Q}
                 \subseteq H^n(\mathcal X_F)^{\mu_N}.
                                                               \tag{H9.107}
\]
In particular neither a new fixed-ambient cross term nor a
new variable algebraic class survives at the geometric
generic parameter.

\paragraph{Actual cycles give the height and radical formula.}
All homologically trivial generic cycles have admissible
closures on these smooth models by H9-T12(iii). For such a
closure $Z$, (H9.107) and actual finite flat cyclic
push--pull show that
\[
                  Z-\frac1NQ^*Q_*Z
\]
is globally homologically trivial. It therefore has zero
intersection with every admissible closure. The height
projection formula gives (H9.102), exactly as in H9-T19.
Surjectivity of corestriction over $\mathbf Q$ proves the
radical equality. For $m=2$, H9-T8 gives the claimed
semipositivity of the old group and hence of the new group.

\paragraph{Freely varying branch points.}
To make the last genericity assertion precise, let $T_B$ be
the open subset of $(\mathbf A^2-\{0\})^{B+2}$ consisting
of ordered homogeneous linear forms
$\ell_0,\ldots,\ell_B,\ell_*$ whose zeros are pairwise
distinct and avoid $\Sigma$. This is a smooth irreducible
parameter space. Form the cyclic root cover associated with
\[
                \ell_0\ell_1\cdots\ell_B\ell_*^{N-1}
                    \in H^0(\mathcal O(B+N)),
                                                               \tag{H9.108}
\]
and normalize. The same explicit local models, with smooth
parameter base $T_B$, prove that this is a smooth proper
family of connected curves with a finite flat degree-$N$
map to $\mathbf P^1\times T_B$. Its pullback pencil family
is smooth and proper as before. The kernel of pushforward
on middle cohomology is the lisse direct summand cut out
by $1-Q^*Q_*/N$.

The family (H9.103) is obtained by setting
$\ell_0=T-sZ$, $\ell_i=T-b_iZ$, and $\ell_*=Z$.
Its kernel has no finite-orbit vector by (H9.104)--(H9.106).
Thus the full $T_B$ kernel has none either, since a
finite-orbit vector remains finite-orbit after restriction
to this subfamily. With constants
$\Omega=\overline{k(T_B)}$, every geometric generic
algebraic kernel class is therefore zero. The actual Chow
push--pull argument again gives (H9.102).

When $N=2$ and $B=2g$, all $2g+2$ exponents in (H9.108)
are one. These are all ordered distinct branch configurations
of a double cover of $\mathbf P^1$; rescaling the linear
forms only changes the defining scalar, which is a square
over $\Omega$. Thus the geometric generic fibre is exactly
the geometric generic hyperelliptic cover of genus $g$.
This verifies the stated scope without a Hurwitz-space
irreducibility theorem or an extension to arbitrary
noncyclic covers.
\end{proof}

\paragraph{What is and is not concluded.}
The theorem is compatible with H9-T17. That example chooses
a special elliptic covering curve having a nonconstant map
to a fixed abelian factor, and produces genuine positive
and negative norm-zero directions. Here the generic
parameter monodromy excludes all new algebraic middle
classes, including the ambient cross term that would
produce those directions. Specialization can create such
classes; no positivity assertion for every special cover
is inferred.

The complete norm kernel is proved height-radical, not
Chow-zero. This is a geometric function-field theorem; it
does not imply an arithmetic Abel--Jacobi injectivity
statement or full positivity of an arithmetic Chow motive.
The only additional background beyond H9-T19 is elementary
tame cyclic-cover normalization and Riemann--Hurwitz, whose
application and all exponents were checked explicitly above.
The middle-convolution source is exactly the directly read
Katz theorem and formulas recorded at the end of H9-T19.
